\section{Zadanie 14}
Z kwadratu jednostkowego \([0, 1] \times [0, 1]\) losujemy punkt w sposób jednostajny.
Później niezależnie wybieramy kąt w sposób jednostajny na \([0, 2\pi]\).
Z jakim prawdopodobieństwem odcinek jednostkowy rozpoczynający się w wylosowanym punkcie
zmierzający w kierunku wylosowanego kąta wychodzi poza rozważany kwadrat?

\subsection*{Rozwiązanie}

\begin{center}
	\begin{tikzpicture}[scale=3.5]
		% Kwadrat [0,1]x[0,1]
		\draw[thick] (0,0) rectangle (1,1);

		% Punkt Q
		\filldraw[red] (0.7,0.4) circle (0.02);
		\node[below left] at (0.7,0.4) {$(x, y)$};

		% Długość odcinka
		\def\L{1}

		% Kąt w stopniach
		\def\kat{160}

		% Punkt R (koniec odcinka)
		\coordinate (R) at ($(0.7,0.4) + (160:\L)$);

		% Odcinek Q -> R
		\draw[blue, thick] (0.7,0.4) -- (R);

		% Zaznaczenie końcowego punktu R
		\filldraw[green] (R) circle (0.02);
		\node[above left] at (R) {$(x+\cos\theta, \; y+\sin\theta)$};

		% Łuk kąta theta (promień 0.25)
		\draw[->, thick, orange]
		(0.7,0.4) ++(0:0.25) arc (0:\kat:0.25) {};

		% Podpis kąta
		\node at ($(0.7,0.4)+(0.02,0.12)$) {$\theta$};
	\end{tikzpicture}
\end{center}

Niech punkt startu to \((x,y)\) i kąt \(\theta\). Koniec odcinka jednostkowego ma współrzędne:
\[
	(x+\cos\theta, \; y+\sin\theta).
\]

Odcinek pozostaje w kwadracie wtedy i tylko wtedy gdy punkt końca odcinka jednostkowego jest w kwadracie, czyli:
\[
	0 \le x + \cos\theta \le 1 \quad \text{oraz} \quad 0 \le y + \sin\theta \le 1.
\]

Dla ustalonego \(\theta\), dopuszczalne wartości \(x\) (z pierwszej nierówności) to:
\[
	x \in [0,1] \cap [-\cos\theta, 1-\cos\theta],
\]
długość tego przedzialu wynosi \(1 - |\cos\theta|\). Analogicznie dla \(y\): \(1 - |\sin\theta|\).

Zatem prawdopodobieństwo, że odcinek pozostaje w kwadracie dla ustalonego \(\theta\), to:
\[
	(1 - |\cos\theta|)(1 - |\sin\theta|).
\]
Teraz wystarczy uśrednić ten wynik po wylosowanym kącie \\
Uśredniając po \(\theta\), otrzymujemy:
\[
	P_{\text{wewnątrz}} = \frac{1}{2\pi} \int_0^{2\pi} (1 - |\cos\theta|)(1 - |\sin\theta|) \, d\theta.
\]

Rozwijamy iloczyn pod całką:
\[
	(1 - |\cos\theta|)(1 - |\sin\theta|) = 1 - |\cos\theta| - |\sin\theta| + |\cos\theta||\sin\theta|.
\]

Zatem:
\[
	P_{\text{wewnątrz}} = 1 - \frac{1}{2\pi} \int_0^{2\pi} (|\cos\theta| + |\sin\theta|) \, d\theta + \frac{1}{2\pi} \int_0^{2\pi} |\cos\theta||\sin\theta| \, d\theta.
\]

Obliczamy całki:
\[
	\int_0^{2\pi} |\cos\theta| \, d\theta = 4, \quad \int_0^{2\pi} |\sin\theta| \, d\theta = 4, \quad \int_0^{2\pi} |\cos\theta||\sin\theta| \, d\theta = 2.
\]

Finalnie:
\[
	P_{\text{wewnątrz}} = 1 - \frac{4+4}{2\pi} + \frac{2}{2\pi} = 1 - \frac{6}{2\pi} = 1 - \frac{3}{\pi}.
\]

Zatem prawdopodobieństwo, że odcinek wychodzi poza kwadrat, wynosi:
\[
	P_{\text{wychodzi}} = \frac{3}{\pi}
\]


\subsection*{Odpowiedź}
\[
	\boxed{\frac{3}{\pi}}
\]
